\documentclass{report}
\usepackage{amsmath,amssymb}
\setlength{\parindent}{0mm}
\setlength{\parskip}{1em}
\begin{document}
\begin{center}
\rule{6in}{1pt} \
{\large Paul Kuberry \\
{\bf An optimization-based approach for interfaces with gaps and overlaps via virtual controls }}

Sandia National Laboratories \\ Mailstop 1320 \\ P O Box 5800 \\ Albuquerque \\ NM 87185-XXXX
\\
{\tt pakuber@sandia.gov}\\
Pavel Bochev\end{center}

Partitioned approaches are popular for solving coupled problems with
possibly different physics in each of one or more subdomains sharing an
interface. Particularly when the physics differ, it may be advantageous
or simply unavoidable to mesh the subdomains independently. In the case
of refining a straight or flat interface, this may result in multiple
discretizations that are coincident but have mismatched vertices. In the
case of refining a curved surface, the resulting discretizations will
likely be noncoincident and mismatched, having gaps and overlaps.

We present an optimization-based approach, which deals with the coupling
between subdomains by introducing a virtual control as a natural boundary
condition on each of the subdomains. Providing boundary conditions to the
physics on each subdomain permits using separate solvers simultaneously.
We introduce terms in the objective of our optimization that penalize any
violation of global flux conservation or violation of continuity between
states. Since the solvers each return a solution on their respective
subdomain, it adds complication to comparing the continuity of states. In
order to measure violations of continuity, we extend solutions between
subdomains by means of extrapolation in a way that makes use of the
gradient. We iteratively update the control until the objective is
minimized. Computational results will be presented that demonstrate
convergence to a manufactured solution, and also that the method can pass
a linear consistency test.


\end{document}
